% Anonymous edition. Compile both documents with python3 build.py.
% Official econsocart class/configuration are redistributed without modification.
\documentclass[ecta,draft]{econsocart}
\ifdefined\pdfminorversion\pdfminorversion=7\fi
\usepackage[utf8]{inputenc}
\usepackage{microtype,mathtools,booktabs,array,enumitem,needspace,longtable,tabularx}
\usepackage[hidelinks,hypertexnames=false]{hyperref}
\usepackage{xurl,bookmark}
% Portable mappings for font shapes used by the official public class.
\DeclareFontShape{T1}{ptm}{m}{uc}{<->ssub*ptm/m/n}{}
\DeclareFontShape{T1}{ptm}{m}{up}{<->ssub*ptm/m/n}{}
\DeclareFontShape{T1}{ptm}{m}{scit}{<->ssub*ptm/m/sc}{}
\hypersetup{pdftitle={The Cost of Voluntary Exchange: Allocation and Settlement under Separate Public Accounts},pdfauthor={},pdfsubject={Anonymous manuscript},pdfkeywords={voluntary exchange; separate public accounts; settlement; reserves}}
\setlength{\emergencystretch}{2em}
\clubpenalty=10000
\widowpenalty=10000
\displaywidowpenalty=10000
\setlist{itemsep=3pt,topsep=5pt}
\setcounter{secnumdepth}{2}
% Neutral masthead for an anonymous working manuscript. No submission is asserted.
\makeatletter
\def\copyright@text{Anonymous Manuscript}
\makeatother
\startlocaldefs
\newtheorem{theorem}{Theorem}
\newtheorem{proposition}[theorem]{Proposition}
\newtheorem{lemma}[theorem]{Lemma}
\newtheorem{corollary}[theorem]{Corollary}
\theoremstyle{definition}
\newtheorem{definition}[theorem]{Definition}
\theoremstyle{definition}
\newtheorem{remark}[theorem]{Remark}
\newtheorem{example}[theorem]{Example}
\newcommand{\R}{\mathbb R}
\newcommand{\FS}{\mathcal F^{S}}
\newcommand{\FP}{\mathcal F^{P}}
\newcommand{\FB}{\mathcal F^{B}}
\newcommand{\argmax}{\operatorname*{arg\,max}}
\newcommand{\argmin}{\operatorname*{arg\,min}}
\newcommand{\supp}{\operatorname{supp}}

\endlocaldefs
\begin{document}
\begin{frontmatter}
\title{The Cost of Voluntary Exchange:\\
Allocation and Settlement under Separate Public Accounts}
\runtitle{The Cost of Voluntary Exchange}
\begin{abstract}
An informed authority chooses allocations and voluntary settlement schedules
under separate public accounts. We characterize implementation losses as
functions of standing capacity, temporary reserves, and approval costs.
Conserving technologies recycle capacity; executable productive releases expand
it. With empty accounts, optimized losses have square-root order under
conservation and linear order with a logarithmic factor under productivity.
Positive standing capacity changes these limiting rates while preserving a
quantitative gap in approval requirements. This gap determines when productive
releases reverse the choice between joint and sequential settlement. The model
retains complete information, divisible irreversible adjustments, balanced
compensation, and unvalued residual public balances.
\end{abstract}
\begin{keyword}
\kwd{Voluntary exchange}
\kwd{public accounts}
\kwd{bilateral clearing}
\kwd{reserves}
\kwd{settlement costs}
\end{keyword}
\end{frontmatter}
\noindent\textsc{JEL classifications:} D47, D61, H41, H55, J26.

\section{Introduction}

An informed public authority chooses an allocation and a voluntary settlement
schedule under separate public accounts. A transaction can consume capacity
before another transaction releases it. This execution constraint changes
the authority's choice between sequential approvals and joint settlement,
even when each institution chooses its own allocation. Conserving transactions
reuse a bounded stock; executable productive releases expand the capacity
available for later approvals. We derive sharp implementation-loss bounds
and identify when this difference reverses the preferred institution.
Balanced compensation secures consent within every approval, while account
feasibility governs its execution. The resulting comparison determines how
standing capacity substitutes for temporary reserves and joint commitment.

The analysis begins with the fully informed planner and its implementation
through unrestricted voluntary trade. We then impose separate-account
feasibility after every bilateral approval. With empty accounts, conserving
technologies incur optimized losses of order $\sqrt{\rho\tau}$, whereas
executable productive releases yield order $\tau\log(\rho/\tau)$, where
$\rho$ prices temporary reserves and $\tau$ prices approvals. Positive
standing balances change both limiting rates to linear order in $\tau$ when
those balances are fixed. The economically relevant difference then lies in
the dependence on standing capacity: inverse under conservation and
logarithmic under productivity. A uniform comparison identifies a range in
which existing balances sustain productive implementation at the minimum-loss
order without rented reserves, while conserving implementation benefits from
additional reserves. The associated welfare gap determines when productive
releases change the preferred settlement institution. Throughout, the informed
authority selects allocations and schedules; adjustments are divisible and
irreversible, and compensation is enforceable and balanced. The model contains
no decentralized bargaining.

\paragraph{Contribution hierarchy.}
The principal result links the resource technology to optimized implementation
losses, jointly accounting for standing capacity, rented reserves, and approval
costs. Theorem~\ref{thm:standingcapacity} establishes the comparison with
positive initial balances; Theorem~\ref{thm:optimizedloss} gives its
empty-account benchmark. The exact reserve--approval frontier identifies the
mechanism, and the general execution conditions extend it beyond the numerical
family. These bounds have two institutional implications. They identify when
existing capacity is sufficient to implement the planner allocation at the optimal loss order
without rented reserves, and when the cost of joint commitment changes the
preferred settlement rule. A common-fee interval establishes the latter
comparison after both institutions choose their allocations.
Corollary~\ref{cor:standingcutoffs} also solves the institutional decision
at the original fixed prices: a strictly positive standing-capacity interval
supports the reversal, and a higher capacity cutoff makes sequential
settlement preferable under both technologies. Table~\ref{tab:contributions} organizes the results around this argument.

\begin{table}[!htbp]
\centering
\caption{Contribution hierarchy}
\label{tab:contributions}
\begin{tabularx}{\linewidth}{@{}>{\raggedright\arraybackslash}p{0.17\linewidth}>{\raggedright\arraybackslash}X>{\raggedright\arraybackslash}p{0.26\linewidth}@{}}
\toprule
Role & Economic content & Results\\
\midrule
Central result & Standing capacity and productive releases determine optimized implementation losses and the value of temporary reserves. & Theorems~\ref{thm:standingcapacity} and~\ref{thm:optimizedloss}\\[3pt]
Exact mechanism & Conservation and productive releases generate inverse and logarithmic approval requirements. & Propositions~\ref{prop:exactproductive} and~\ref{prop:standingexact}; Theorem~\ref{thm:generalproductive}\\[3pt]
Institutional implication & Commitment thresholds identify an endogenous-allocation reversal at common prices. & Corollaries~\ref{cor:standingreversal} and~\ref{cor:standingcutoffs}; Theorem~\ref{thm:endogenousreversal}\\[3pt]
Economic margins & Allocation-dependent fees, distorted quantities, binding participation, and discrete decisions. & Propositions~A.5%<xref:prop:allocationfees>
 and~A.6%<xref:prop:allocationdistortion>
; Corollary~\ref{cor:gridrobustreversal}\\[3pt]
Foundation & Balanced implementation and exact recovery with a repayable reserve. & Section~\ref{sec:trade}; Theorem~A.2%<xref:thm:temporaryreserve>
\\[3pt]
Application & Priority changes the number and identity of water trade requests served by a fixed capacity. & Section~\ref{sec:waterapplication}; Appendix~M%<xref:app:waterbenchmark>
\\
\bottomrule
\end{tabularx}
\end{table}

\paragraph{The approval frontier and its boundary.}
A blocked start and a strictly positive nonexpanding weighted stock imply
inverse-reserve approval requirements. An executable productive sequence
supports geometric expansion and logarithmic requirements when its releases
also permit completion at the planner allocation. We state both sets of
conditions directly and provide a finite family of linear programs that
constructs the productive sequence from the resource matrix, preferences,
caps, and planner value. Section~\ref{sec:standingcapacity} extends the
comparison uniformly to positive standing balances. The resulting bounds
separate three policy margins: using standing capacity, authorizing
temporary reserves, and paying for joint commitment. The exact constants
and parity terms in the three-account family sharpen this general comparison.

\paragraph{Mandate and extensions.}
The core sets $B\equiv0$: employee surplus is maximized subject to separate
public protections, with residual balances carrying no additional welfare
value. Implementation costs are financed from separately authorized operating
resources and deducted from aggregate surplus. A curved-benefit example shows
how sequential costs reduce the selected quantity; a participation example
shows how employee consent raises the reserve required by a feasible physical
schedule. The institutional reversal also holds over a stated interval of
allocation-dependent fees and for a fixed adjustment grid with positive
initial balances. Appendix~K%<xref:app:valuedbalances>
 studies valued public
balances and the associated financing conditions.
Appendix~L%<xref:app:indivisibility>
 derives discrete reserve floors and
convergence to the divisible frontier.

\paragraph{Administrative application.}
Section~\ref{sec:waterapplication} examines a published water trade queue
with positive initial capacity. The same 203 submitted requests produce 98
successful applications under recorded random priority, 106 under arrival
order, and a maximum of 159 under smallest volume first. Executed water
volume changes by less than $0.1\%$. The application isolates the allocation
consequences of priority conditional on the submitted requests and other
account requirements. The exact count benchmark and administrative replay
give an observable institutional example of why high aggregate capacity use
can coexist with substantial differences in access across applications.

\paragraph{Relation to existing work.}
The bilateral-exchange literature characterizes the relationship between
pairwise improvements and efficiency
\citep{Feldman1973,Madden1978,GoldmanStarr1982,Fisher1989}.
We study the execution cost of voluntary improvements when separate public
accounts constrain every intermediate allocation. Continuous-net theory
provides finite-reachability characterizations and divided-execution
constructions \citep{JulvezRecaldeSilva2003,FracaHaddad2015}.
Our economic comparison optimizes the allocation and its voluntary
implementation together. Preferences and balanced compensation determine
admissible atomic approvals; account protections determine their executable
order. The reserve frontier prices this interaction and yields thresholds
for institutional choice, including with positive standing capacity.
Decomposing an atomic approval into separate activity moves can change
consent, resource requirements, and implementation cost.

Settlement economics establishes that the organization of execution affects
resource use and welfare. \citet{KahnRoberds2001} derive the cost of gross
settlement from credit constraints and analyze the role of intraday credit
and collateral. \citet{KahnMcAndrewsRoberds2003} connect the underlying trading
economy to gridlock, collateral requirements, and settlement delay.
\citet{MartinMcAndrews2008} study how queue design changes incentives and
welfare, while \citet{AtalayMartinMcAndrews2010} compare settlement outcomes
with the planner's allocation. Related work studies settlement institutions
in trading economies \citep{KoepplMonnetTemzelides2008}.
Our mechanism concerns how voluntary execution replenishes the particular
public capacities needed for subsequent approvals. Its economic consequence
is a change in optimized implementation losses and in the authority's
willingness to pay for joint authorization. The conserving cost balance
follows the inventory logic of \citet{Baumol1952}; productive releases change
the rate at which additional capacity reduces the approval burden.
The authority chooses the allocation and schedule under complete information.
Private-information and computational-complexity questions are developed
separately.

\section{The planner and unrestricted trade}
\label{sec:planner}

\subsection{Preferences, accounts, and the settlement mandate}

There are finite sets of demanders $D$ and suppliers $S$. Demander $i$ chooses
$a_i\in[0,\bar a_i]$ and receives benefit $U_i(a_i)$; supplier $j$ chooses
$b_j\in[0,\bar b_j]$ and incurs cost $C_j(b_j)$. Roles are fixed and publicly
verified. Let $z=(a,b)$ and let $X$ be the product of these intervals, with
positive finite caps. The baseline is $z=0$. Benefits and costs are continuously
differentiable on neighborhoods of their intervals, benefits are concave,
and costs are convex and nondecreasing. Normalize
$U_i(0)=C_j(0)=0$. The frontier and cost results specialize to positive linear
values and costs, $U_i(a_i)=v_i a_i$ and $C_j(b_j)=c_jb_j$.

There are $K$ public resource accounts. Their initial capacities are
$R\geq0$, and $M$ is the verified linear resource-effect matrix. A positive
entry of $Mz$ releases capacity; a negative entry consumes it. The feasible set is
\begin{equation}
 \FS(R)=\{z\in X:R+Mz\geq0\}.
 \label{eq:feasible}
\end{equation}
Inequalities between vectors are coordinatewise. Rows can distinguish dates,
grades, or organizational units whose authorizations cannot be combined. For
one demander, the retirement interpretation has pension, payroll, and position
accounts:
\begin{equation}
R+Mz=\begin{pmatrix}
 R^P-\kappa_da+\sum_j\kappa_jb_j\\
 R^W+\omega_da-\sum_j\omega_jb_j\\
 R^H+a-\sum_jb_j
\end{pmatrix},
\qquad \kappa_d,\omega_d,\kappa_j,\omega_j>0.
\label{eq:threeaccounts}
\end{equation}
The pension row is a specified ceiling on incremental pension exposure, not an
agency-held retirement fund. Payroll effects must be recognized within the
relevant authorization, and position units represent compatible service time.
The institutional technology grants authority to certify and reuse these releases.

Compensation $t_h$ is money received by employee $h$. Employee utilities are
$U_i(a_i)+t_i$ and $-C_j(b_j)+t_j$. Compensation is exactly balanced when
$\sum_h t_h=0$; these transfers do not enter $M$. They do not amend an
appropriation or create a position. Define employee surplus by
\begin{equation}
 W(z)=\sum_{i\in D}U_i(a_i)-\sum_{j\in S}C_j(b_j)
     =\sum_h f_h(z_h).
 \label{eq:welfare}
\end{equation}

\begin{center}
\fbox{\begin{minipage}{0.92\linewidth}
\textbf{Settlement mandate.} The core sets $B\equiv0$: the authority maximizes
employee surplus subject to every separate public-account constraint, with no
additional welfare weight on residual balances and no net public payment to
employees. Information is complete, adjustments are divisible and irreversible,
and compensation is enforceable and balanced. Every account remains nonnegative
after each approval. Temporary reserves must be returned; original balances may
be consumed. Implementation expenditures come from separately authorized
operating resources.
\end{minipage}}
\end{center}

The fully informed benchmark under this mandate is
\begin{equation}
 V^S(R)=\max_{z\in\FS(R)}W(z).
 \tag{SP}\label{eq:SP}
\end{equation}
Compactness and continuity give an optimum; the baseline makes its value
nonnegative. This mandate fixes public protections as constraints. A broader
appraisal that values residual capacity is developed in
Appendix~K%<xref:app:valuedbalances>
.

\subsection{The endogenous planner allocation in the productive family}
\label{sec:solvedplanner}
\label{sec:localrobustness}

The central family has one demander, two suppliers, unit caps, zero initial
balances, $v=1$, $c_1=c_2=1/10$, and
\[
 M_h=\begin{pmatrix}
 -2&3/2+h&5/2\\
 1&-1/2&-3/2\\
 1&-1&-1
 \end{pmatrix},\qquad 0\leq h<1/2.
\]
Only supplier 1's pension release changes with $h$. Preferences, caps, the
other resource effects, and bilateral permissions stay fixed.

\begin{proposition}[The endogenous planner optimum]
\label{prop:productiveplanner}
For each demander quantity $a\in[0,1]$, the minimum supplier cost under the
separate-account constraints is
\[
 C_h^S(a)=\frac{1+h}{5(2+3h)}a,
 \qquad
 (b_1,b_2)=\left(\frac{a}{2+3h},
                     \frac{(1+2h)a}{2+3h}\right).
\]
The planner optimum and value are therefore
\[
 z_h^*=\left(1,\frac1{2+3h},\frac{1+2h}{2+3h}\right),
 \qquad V_h^*=1-\frac{1+h}{5(2+3h)}.
\]
The optimum is unique, and its residual balances are
$M_hz_h^*=(0,0,h/(2+3h))^{\mathsf T}$.
\end{proposition}

\begin{proof}
At a fixed $a$, pension and payroll feasibility imply
\[
 (3/2+h)b_1+(5/2)b_2\geq2a,
 \qquad -(1/2)b_1-(3/2)b_2\geq-a.
\]
Multiplying these inequalities by $2/(2+3h)$ and
$2(1-h)/(2+3h)$, respectively, and adding gives
\[
 b_1+b_2\geq\frac{2(1+h)}{2+3h}a.
\]
The stated bundle attains this bound and leaves position capacity
$ha/(2+3h)\geq0$. Its quantities respect the unit caps. Both multipliers
are positive, so equality requires both fiscal constraints to bind; their
independent equations give the unique cost-minimizing bundle. Employee surplus
at this bundle is $a[1-(1+h)/(5(2+3h))]$, whose coefficient is positive.
Hence the unique optimum has $a=1$.
\end{proof}

At $h=1/4$, the optimum is $(1,4/11,6/11)$ and its value is $10/11$.
The proportional target $(1,1/2,1/2)$ has value $9/10$ and leaves pension
capacity $1/8$. It is feasible but not planner-optimal: it uses more supplier
activity to obtain the same demander benefit. In general,
\[
 V_h^*-\frac9{10}=\frac{h}{10(2+3h)},
\]
which is strictly positive for $h>0$ and equals $1/110$ at $h=1/4$.

\subsection{Unrestricted trade implements the benchmark}
\label{sec:trade}

An unrestricted agreement makes an entire feasible allocation and its
compensation effective together. Each employee may retain the baseline, and
there are no wealth constraints or transfer caps. At any allocation with
$W(z)\geq0$, choose $\alpha_h\geq0$ summing to one and set
$t_h=\alpha_hW(z)-f_h(z_h)$. Transfers balance and give each employee the
nonnegative payoff $\alpha_hW(z)$. Thus unrestricted voluntary agreement
implements a planner optimum.

The allocation also has competitive support. At account prices $p\geq0$,
employee $h$ chooses $z_h$ to maximize $f_h(z_h)+p^{\mathsf T}m_hz_h$,
where $m_h$ is its column of $M$. Concavity and the polyhedral normal-cone
condition at an optimum supply supporting prices with
$p^{\mathsf T}(R+Mz^*)=0$. Fixed nonnegative rebates totaling
$p^{\mathsf T}R$ make employee transfers exactly balanced. Each employee
weakly prefers its choice to zero. Conversely, summing the maximizing
inequalities establishes planner optimality. This standard welfare argument
\citep{ArrowDebreu1954,BoydVandenberghe2004} is proved, including valued public
balances, in Proposition~A.8%<xref:prop:publiccompetitive>
 of
Appendix~K%<xref:sec:implementation>
.

For $M_h$, prices
\[
 p_h=\left(\frac1{5(2+3h)},
                 \frac{1-h}{5(2+3h)},0\right)
\]
pay each supplier $1/10$ per unit and charge the demander
$(1+h)/[5(2+3h)]$ per unit. The demander chooses one; the cost-minimizing
supplier quantities clear the accounts and satisfy supplier optimality.
The corresponding transfers leave the demander $V_h^*$ and both suppliers
zero. These prices support the planner allocation through jointly effective
agreements under complete information.

\section{Constrained settlement and the exact reserve frontier}
\label{sec:recovery}\label{sec:institution}
\label{sec:sequential}\label{sec:sharpseeds}\label{sec:seeddesign}

A joint package takes effect when its aggregate allocation satisfies the ledger.
A sequential approval changes at most one demander and one supplier activity;
demander-only approvals are permitted. Activities only increase. Each approval
must leave every account nonnegative and give each participating employee a
nonnegative incremental gain under exactly balanced compensation. With linear
preferences, increments $(x,y)$ for pair $(i,j)$ satisfy this participation
condition precisely when $v_ix\geq c_jy$. Previously settled releases can be
reused. An independently cleared pair cannot use another pair's release; the
sequential institution studied here expressly permits such reuse.

Let $V_T^{\mathrm{seq}}(R)$ denote the largest employee surplus attainable
from $R$ in at most $T$ such approvals. This notation alone imposes no repayment
condition. In every institutional comparison below, a reserve $\epsilon r$
added to an original ledger $R^0$ must be returned, so the terminal allocation
also satisfies $R^0+Mz\geq0$. Reserve authorization and return are not bilateral
approvals. All costs below price this same convention.

Two implementation facts motivate the frontier. First, no nonzero bilateral
approval is feasible from zero in $M_h$: supplier 1 requires $x\geq y$ from
positions but $x\leq(3/4+h/2)y<y$ from pensions; supplier 2 requires
$x\geq3y/2$ from payroll but $x\leq5y/4$ from pensions. A demander-only move
also consumes pension capacity. Second, a reserve positive in every account
permits exact recovery of an original positive-surplus planner optimum under
nondecreasing concave benefits, convex costs with positive derivatives, and
complete bilateral access. Theorem~A.2%<xref:thm:temporaryreserve>
 proves this
claim by decomposing the optimum into employee-voluntary bundles and executing
sufficiently small rounds. The reserve covers within-round shortfalls and is
returned at completion. Its full statement and proof are in
Appendix~A%<xref:app:designproofs>
.

For $\theta\in(0,1]$, define the proportional target
$z_\theta=(\theta,\theta/2,\theta/2)$ and let
$N_h(\theta,\epsilon)$ be its minimum number of voluntary approvals from
$\epsilon(1,1,1)$. Its surplus is $9\theta/10$ and its net release is
$(h\theta/2,0,0)$, so the reserve is repayable. This target provides an exactly
solvable intermediate stage. Proposition~\ref{prop:productiveplanner} shows
why it is not the planner optimum when $h>0$.

\begin{proposition}[The exact reserve frontier with productive releases]
\label{prop:exactproductive}
In the economy $M_h$, fix $\theta\in(0,1]$ and an integer $T\geq2$.
The least reserve per account that permits exact voluntary implementation
of $z_\theta$ in at most $T$ bilateral approvals is
\begin{equation}
 \epsilon_T(h,\theta)=
 \begin{cases}
 \displaystyle\frac{h\theta}{6(G_T(h)-1)},&0<h<1/2,\\[6pt]
 \displaystyle\frac{\theta}{12(T-1)},&h=0,
 \end{cases}
 \label{eq:exactproductiveseed}
\end{equation}
Here $\gamma_h=(1+2h)/(1-2h)$ and, for $n\geq1$,
\[
 G_{2n}(h)=\frac{1+\gamma_h}{2}\gamma_h^{\,n-1},
 \qquad G_{2n+1}(h)=\gamma_h^{\,n}.
\]
An attaining schedule alternates suppliers, beginning with supplier 1.
Every intermediate account balance is nonnegative, each transaction is
voluntary under balanced compensation, and the reserve can be returned
after the final approval. The formula for $h=0$ is also the continuous
limit as $h\downarrow0$ for each fixed $T$ and $\theta$.
\end{proposition}

\begin{proof}[Proof sketch]
A change of variables projects feasible balances into a two-dimensional
conversion problem. In that relaxation, deleting moves without supplier
activity and combining adjacent moves using the same supplier leaves an
alternating word. Successive maximal conversions move between its boundary
faces; the associated linear recurrence
has multiplier $\gamma_h$ and the displayed parity term $G_T$. Bounding all
alternating words by this recurrence gives necessity for physical schedules.
Lifting the boundary construction back to activity increments attains the bound,
has positive surplus at every approval,
respects the activity caps, and returns the reserve. At $h=0$ the recurrence
becomes additive. Appendix~B%<xref:app:localrobustness>
 proves both bounds and
verifies the attaining transfers and terminal balances.
\end{proof}

The frontier has inverse-reserve order at $h=0$ and logarithmic order at every
fixed $h>0$. Conservation keeps the usable weighted
stock proportional to the initial reserve. Productive rounds create capacity
that later rounds reuse. Section~\ref{sec:generalproductive} gives conditions
for this distinction beyond $M_h$ and for completion at an endogenous planner
optimum.

\section{Productive releases and optimized implementation losses}
The authority authorizes reserve $\epsilon r$ at cost $\rho\epsilon$, with
$\rho>0$, and pays $\tau>0$ for each approval. A joint package uses one
approval and pays a fixed contingent-commitment cost $F$. Unless otherwise
specified, $F$ is the same for every nonzero joint package, and both institutions
have common verification cost $\mathcal K=0$. The zero allocation costs nothing. Costs
are paid from preauthorized operating resources and deducted from aggregate
surplus; employee compensation remains balanced within each approval. If these
costs consume a constrained account, that consumption must instead be included
in the ledger. No exogenous reserve cap is imposed here.

\paragraph{Authorization capacity and the institutional comparison.}
An account balance authorizes specified adjustments within one planning
episode. A reserve temporarily enlarges that authorization and is returned
at the terminal allocation. Its price $\rho\epsilon$ measures the opportunity
or administrative cost of making the reserve vector available. The retirement
application interprets these balances as permissions to recognize and reuse
pension, payroll, and position releases. The reserve composition $r$ is fixed.

We model two authorization technologies. Under routine approval, a bilateral
order becomes effective when issued, compensation is enforceable, and each
account officer checks the resulting ledger against the capacity then
available. The authority can plan the entire sequence in advance. A planned
future release becomes usable when the order producing it takes effect.
Under joint authorization, the officers certify a common program whose
debits and releases become effective together. The program requires a
contingent authorization from each protected account and a common
effectiveness instruction. Let $f_k\geq0$ be the resource cost of preparing
that contingent authorization for account $k$, and let $f_0\geq0$ be the cost
of activating the common instruction. With fixed administrative coverage,
the additional program cost is
\[
 F=f_0+\sum_{k\in\mathcal A}f_k,
\]
where $\mathcal A$ is the fixed set of protected accounts.
The ordinary approval cost $\tau$ covers the final authorization under either
technology. The costs $f_k$ purchase conditional effectiveness across orders;
the compensation contracts within each approved order are enforceable under
both technologies. This administrative technology therefore separates the
ability to plan a program from the authority to make its component orders
jointly effective. Complete information removes uncertainty about gains and
resource effects; it leaves these authorization requirements in place.
In financial settlement, linking the completion of
obligations is likewise an explicit design requirement
\citep[Principle~12]{CPSSIOSCO2012}.

The cost $F$ is an avoidable resource cost of the program under consideration.
When the required joint facility is already available at zero incremental
cost, the relevant value is $F=0$. When administrative coverage depends on
the chosen program, the same construction yields $F(z)$, as studied in
Proposition~A.5%<xref:prop:allocationfees>
. These costs are evaluated at common
administrative prices across the technologies being compared. With a common
fixed charge, a schedule containing joint approvals is weakly dominated by
one joint package implementing its terminal allocation: balanced compensation
supports that endpoint with one approval and no rented reserve. The two
settlement rules thus exhaust the useful schedules within this authorization
technology. Variable charges permit additional batching choices when those
choices are included in the institutional menu.

\subsection{Executable productive releases and the boundary}
\label{sec:generalproductive}

The geometric construction has a sufficient condition that does not
depend on the coefficients of $M_h$. We state it for a finite linear
economy with zero initial balances, positive values and costs, finite
action caps, and reserve $\epsilon r$, where $r\gg0$. Retain complete
information, irreversible divisible adjustments, balanced compensation,
and the $B\equiv0$ mandate. Every demander--supplier pair and
demander-only move is permitted. A voluntary bundle has at most one demander
and one supplier coordinate, or a demander coordinate alone; a pair
$(x,y)$ satisfies $v_i x\geq c_jy$. Assume a blocked start: no nonzero
voluntary bundle is feasible from zero. These assumptions define the quantitative class; the recovery theorem in
the appendix also covers concave benefits and convex costs. Let $\mathcal D$ be
the union of these voluntary directions normalized to have coordinate sum one,
and put $\bar v=\max_i v_i$.

\begin{definition}[An executable productive release]
\label{def:executableproductive}
A feasible positive-surplus target $z^*$ admits an executable productive
release if there are $\alpha,\kappa>0$, nonzero voluntary bundles
$u^1,\ldots,u^k$, and a possibly empty finite list of nonzero voluntary
bundles $v^1,\ldots,v^\ell$, such that, writing
\[
 d=\sum_{j=1}^k u^j,\qquad p=Md,\qquad
 c=\sum_{j=1}^\ell v^j,
\]
one has
\[
 p\geq0,\quad p\ne0,\quad z^*=\alpha d+c,
 \qquad
 M\sum_{j=1}^s u^j\geq-\kappa p
 \quad(0\leq s\leq k),
\]
and
\[
 \alpha p+M\sum_{j=1}^s v^j\geq0
 \quad(0\leq s\leq\ell).
\]
All inequalities between resource vectors are coordinatewise.
\end{definition}

The vector $p$ is the net capacity released by one productive bundle.
The prefix inequalities require that this same capacity can finance the
bundle's temporary shortfalls. The final list then completes the target
using the accumulated release. The prefix and completion inequalities distinguish an executable gain
from a net release that cannot finance its own implementation. No
positivity is required in every coordinate of $p$; resources with
$p_i=0$ must be generated before use within the specified bundle.
Since all activities are nonnegative and the completed target respects
the caps, every partial implementation also respects them.

\begin{theorem}[A general logarithmic frontier and its boundary]
\label{thm:generalproductive}
Suppose the preceding assumptions hold and $z^*$ admits an executable
productive release. Put
\[
 \sigma=\min_{i:p_i>0}\frac{r_i}{p_i}>0,
 \qquad \Lambda=1+\frac1\kappa>1.
\]
The minimum number $N^*(\epsilon)$ of voluntary approvals implementing
$z^*$ with full reserve repayment satisfies
\[
 N^*(\epsilon)
 \leq k\left\lceil
 \frac{\log(1+\alpha/(\sigma\epsilon))}{\log\Lambda}
 \right\rceil+\ell,
 \qquad
 N^*(\epsilon)=\Theta(\log(1/\epsilon))
 \quad\text{as }\epsilon\downarrow0.
\]
In contrast, in a blocked-start economy possessing strictly positive
weights $\lambda$ with
\[
 \lambda^{\mathsf T}Md\leq0
 \quad\text{for every normalized voluntary direction }d\in\mathcal D,
\]
implementing any positive-surplus planner optimum with a fully repayable
reserve requires $\Theta(1/\epsilon)$ approvals.
\end{theorem}

\begin{proof}[Proof sketch]
From stock $s p$, the prefix inequalities permit a scaled productive
bundle of size $s/\kappa$. Its release raises that stock to
$\Lambda s p$. Starting with $\epsilon r\geq\sigma\epsilon p$,
repeat this operation and truncate the last repetition when cumulative
bundle size equals $\alpha$. Execute the completion list and return the
reserve. The blocked-start bound in
Theorem~A.1%<xref:thm:generalseedbounds>
 supplies the matching logarithmic
lower bound. Under nonexpansion, the weighted stock remains proportional
to $\epsilon$, giving the inverse-reserve lower bound; divided execution
gives the matching upper bound. Appendix~C%<xref:app:generalproductive>

provides the complete argument and verifies membership of every
$M_h$, $h\in(0,1/2)$.
\end{proof}

The nonexpansion and productive-execution conditions identify two distinct
technological regimes. The former bounds the weighted stock available to
each approval; the latter supplies an executable sequence that expands it
and completes the planner allocation. Technologies outside these sufficient
classes require a separate execution analysis. The coefficients of $M_h$
determine its exact frontier, while the general loss bounds depend on the
resource and execution conditions.

\subsection{A sufficient test stated in technological primitives}
\label{sec:primitiveproductivity}

The executable-release condition can be certified without supplying a
schedule in advance. The following finite collection of linear programs
uses only $M$, values, costs, and action caps. It constructs a sufficient certificate for logarithmic implementation. Maintain the
zero baseline, blocked start, positive linear preferences, complete
bilateral access, and $r\gg0$ of Section~\ref{sec:generalproductive}.
First solve the planner's linear program and denote its value by $V^*>0$.

Let $\mathcal T=(D\times S)\cup\{(i,0):i\in D\}$ be the finite set
of bilateral and demander-only transaction types, and let
$L=|D||S|+|D|$. For a pair $(i,j)$, define the cone
\[
 \mathcal C_{ij}
 =\{x e_i+y e_j:x,y\geq0,\ v_ix\geq c_jy\};
 \qquad \mathcal C_{i0}=\{x e_i:x\geq0\}.
\]
Fix a nonempty subset $R$ of resource accounts and a permutation
$\pi$ of $\mathcal T$. Solve
\[
 \begin{aligned}
 L_{R,\pi}:={}&\max \eta\\
 \text{subject to }&u^s\in\mathcal C_{\pi(s)}\quad(s=1,\ldots,L),
       \qquad z=\sum_{s=1}^L u^s,\\
 &0\leq z\leq\bar z,\qquad W(z)=V^*,\qquad\eta\geq0,\\
 &(Mz)_k\geq\eta\quad(k\in R),
       \qquad(Mz)_k=0\quad(k\notin R),\\
 &\left(M\sum_{s=1}^t u^s\right)_k\geq0
       \quad(k\notin R,\ t=1,\ldots,L).
 \end{aligned}
\]
Infeasible programs fail the test. The constraints permit a type to
receive zero activity. They require the program to choose a planner
optimum and divide it among voluntary transaction types. Accounts
outside $R$ must be replenished before use; temporary shortfalls in
$R$ can subsequently be funded by the positive release in those
same accounts. This is a finite test on primitives, although the
enumeration of supports and orders need not be computationally cheap.

\begin{proposition}[A primitive certificate for logarithmic implementation]
\label{prop:primitiveproductivity}
If $L_{R,\pi}>0$ for at least one support and ordering, an original
planner optimum has a voluntary implementation with full reserve
repayment and $\Theta(\log(1/\epsilon))$ approvals. More explicitly,
take an optimal solution of that program, delete its zero bundles, and
write $\nu$ for the number remaining and $p=Mz$. Define
\[
 \kappa=
 \max_{k\in R,\,1\leq t\leq L}
 \frac{[-(M\sum_{s=1}^t u^s)_k]_+}{p_k},
 \qquad
 \sigma=\min_{k\in R}\frac{r_k}{p_k}.
\]
Then $\kappa,\sigma>0$, and the optimum is implementable with at most
\[
 \nu\left\lceil
 \frac{\log(1+1/(\sigma\epsilon))}{\log(1+1/\kappa)}
 \right\rceil
\]
approvals. If a technology $M_C\leq M$ with the same preferences,
caps, and trading permissions has a blocked start, positive planner
value, and a strictly positive nonexpansion certificate, then
Corollary~\ref{cor:optimizedreversal} applies. Thus the primitive test
also supplies a nonempty interval of common fixed commitment
fees at sufficiently small approval prices at which the optimized institutional ranking reverses.
\end{proposition}

\begin{proof}[Proof sketch]
The program itself supplies a planner optimum and its voluntary
decomposition. Its positive value makes $p$ strictly positive exactly
on $R$. The prefix inequalities outside $R$ and the definition of
$\kappa$ give every prefix the lower bound $-\kappa p$. A blocked
start rules out $\kappa=0$. Hence the output satisfies
Definition~\ref{def:executableproductive} with $d=z$, $\alpha=1$,
and no completion list. Apply Theorem~\ref{thm:generalproductive}
and then Corollary~\ref{cor:optimizedreversal}.
Appendix~D%<xref:app:primitiveproductivity>
 supplies details and verifies
that the test passes for every $M_h$, $h\in(0,1/2)$.
\end{proof}

A successful program constructs both the planner endpoint and its
productive path. Allowing repeated transaction types or a separate completion
stage enlarges the class of certificates beyond this finite test.

\subsection{An interacting three-stage release cycle}
\label{sec:interactingexample}

The test also covers a technology with four accounts and three
suppliers linked through one demander:
\[
 \widetilde M_h=
 \begin{pmatrix}
 -1&1&1&1\\
 1&-2&-1&h\\
 0&1&-1&0\\
 0&0&1&-1
 \end{pmatrix},\qquad h\geq0.
\]
Let the demander value be one, each supplier cost be $1/10$, the
demander cap be one, and each supplier cap be $1/3$. The first
account implies $b_1+b_2+b_3\geq a$, so the unique planner optimum
and its value are
\[
 \widetilde z^*=(1,1/3,1/3,1/3),\qquad
 \widetilde V^*=9/10,
 \qquad
 \widetilde M_h\widetilde z^*=(0,h/3,0,0)^{\mathsf T}.
\]
At this optimum every supplier cap binds. Each bilateral first move
is blocked, including when $h>0$. Three voluntary pairs with equal
demander and supplier increments transfer usable capacity successively
from account 2 to account 3, from 3 to 4, and from 4 to 2; the last
transfer returns $1+h$ units for each unit it consumes. The LP test
passes with $R=\{2\}$ and order $(1,2,3)$ of suppliers. Its
constants are $\kappa=1/h$ and $\sigma=3/h$ for the uniform
reserve direction. This is an interacting resource cycle, with a
shared demander and three suppliers, rather than a block-diagonal
copy of the three-account example.

For a concrete optimized comparison, take $h=0$ and $h=1$, reserve
price $\rho=1$, approval cost $\tau=10^{-6}$, and fixed joint fee
$F=1/5000$. The best joint value in either economy is $0.899799$.
Under conservation the nonexpansion bound gives
\[
 \mathcal V_0^{\mathrm{seq}}
 \leq\frac9{10}-\sqrt{\frac3{10^7}}
 <0.899799.
\]
At $h=1$, reserve $10^{-6}$ and at most $57$ approvals implement the
planner optimum with full repayment, giving
\[
 \mathcal V_1^{\mathrm{seq}}
 \geq\frac{449971}{500000}=0.899942>0.899799.
\]
The allocations are freely chosen in both comparisons. The first
inequality is a global upper bound and the second an attainable lower
bound, which together establish the reversal without solving either
sequential optimization problem exactly. All values, costs, and caps
are unchanged across technologies. The complete primitive calculations
are in Appendix~D%<xref:app:primitiveproductivity>
.

\subsection{Optimized welfare losses and institutional thresholds}
\label{sec:optimizedloss}

Institutional welfare depends on the allocation as well as its execution.
We measure the loss after optimizing jointly over allocations, reserves,
and schedules.
For a zero-balance economy with planner value $V_M>0$, define
\[
 \begin{aligned}
 \mathcal V_M^{\mathrm{seq}}(\rho,\tau)
 &=\max\left\{0,\sup_{(z,\epsilon,T)\ \mathrm{admissible}}
       [W(z)-\rho\epsilon-\tau T]\right\},\\
 L_M(\rho,\tau)&=V_M-\mathcal V_M^{\mathrm{seq}}(\rho,\tau).
 \end{aligned}
\]
Admissibility includes employee-voluntary bilateral approvals, all action
caps and intermediate account constraints, and full repayment of
$\epsilon r$. The baseline incurs no implementation cost. In particular,
$0\leq L_M\leq V_M$; the loss includes any reduction in gross surplus
chosen to economize on implementation.

\begin{theorem}[Optimized implementation losses]
\label{thm:optimizedloss}
Consider the positive-linear, complete-access class of
Section~\ref{sec:generalproductive}, with zero original balances,
$r\gg0$, a blocked start, and $V_M>0$. Fix $\rho>0$.
As $\tau\downarrow0$:
\begin{enumerate}
 \item If the technology admits a strictly positive nonexpansion
 certificate, then
 \[
       L_M(\rho,\tau)=\Theta(\sqrt{\rho\tau}).
 \]
 \item If a planner optimum admits an executable productive release
 in Definition~\ref{def:executableproductive}, then
 \[
       L_M(\rho,\tau)=\Theta\bigl(\tau\log(\rho/\tau)\bigr).
 \]
\end{enumerate}
The constants depend on the fixed economy, its reserve direction, and
the stated certificates. Both rates concern the optimized sequential
value over all allocations, not merely the cost of a particular planner
implementation. For all sufficiently small $\tau$, that optimized value
is strictly positive.
\end{theorem}

\begin{proof}[Proof sketch]
In the nonexpanding class, any implemented surplus $W>0$ requires
$T\geq aW/\epsilon$, where
$a=\delta_\lambda/(\bar vS_\lambda)>0$. Minimizing the resulting
reserve and approval costs gives $2\sqrt{\rho\tau aW}$.
The function $W-2\sqrt{\rho\tau aW}$ is convex on $[0,V_M]$;
its maximum occurs at an endpoint. Thus allowing the institution to
reduce its allocation cannot remove the square-root lower bound.
Divided execution of a planner optimum attains the matching order.

In the productive class, an allocation with $W<V_M/2$ already loses
at least $V_M/2$ before implementation costs. For $W\geq V_M/2$,
the blocked-start bound gives
$T\geq\beta^{-1}\log(1+c/\epsilon)$ for constants
$\beta,c>0$. Minimizing
$\rho\epsilon+(\tau/\beta)\log(1+c/\epsilon)$ gives a lower
bound of logarithmic order. Executing the productive certificate at
$\epsilon=\tau/\rho$ gives the matching upper bound. These feasible
planner implementations have costs tending to zero, establishing
positive adoption. Appendix~E%<xref:app:optimizedloss>
 gives the bounds
and the full proof.
\end{proof}

Optimized implementation loss determines the authority's willingness to pay
for a common effectiveness instruction. Joint authorization saves the
sequential loss but still requires one ordinary approval. The remaining
saving is the maximum additional program cost the authority will pay.

\begin{proposition}[The joint-commitment threshold]
\label{prop:feethreshold}
Let every nonzero joint package cost $\tau+F$, where $F\geq0$,
and keep the sequential costs above. If
$\mathcal V_M^{\mathrm{seq}}(\rho,\tau)>0$, define
\[
 F_M^*(\rho,\tau)=L_M(\rho,\tau)-\tau.
\]
Joint settlement has higher optimized net surplus when $F<F_M^*$;
the institutions tie when $F=F_M^*$; sequential settlement has higher
optimized net surplus when $F>F_M^*$. Under the blocked-start assumption,
$F_M^*\geq\tau>0$.

With $F=0$, joint settlement weakly dominates sequential settlement
for every $\rho,\tau>0$. In the blocked-start class the dominance is
strict whenever $V_M>\tau$. If instead
$\mathcal V_M^{\mathrm{seq}}=0$, joint settlement is strictly preferred
exactly when $F<V_M-\tau$; otherwise both institutions select the baseline.
\end{proposition}

\begin{proof}[Proof sketch]
The joint value is $\max\{0,V_M-\tau-F\}$, whereas the
sequential value is $V_M-L_M$. Comparing these expressions proves
the threshold statements, including the baseline qualification.
Every nonzero fully repayable schedule in a blocked-start economy
requires at least two approvals. Hence
$\mathcal V_M^{\mathrm{seq}}\leq\max\{0,V_M-2\tau\}$ and,
when adoption is positive, $L_M\geq2\tau$. More generally, a joint
package with $F=0$ implements any sequential endpoint with one approval
and no reserve cost. Appendix~E%<xref:app:optimizedloss>
 records the argument.
\end{proof}

\begin{corollary}[The interval supporting an endogenous reversal]
\label{cor:optimizedreversal}\label{thm:generalreversal}
Compare a nonexpanding economy $C$ and an executable productive economy
$P$ satisfying the respective hypotheses of
Theorem~\ref{thm:optimizedloss}, with common preferences, action caps,
reserve direction, bilateral permissions, and $\rho>0$. For all
sufficiently small $\tau>0$, both sequential values are positive and
\[
 F_C^*=\Theta(\sqrt{\rho\tau}),\qquad
 F_P^*=\Theta\bigl(\tau\log(\rho/\tau)\bigr),
 \qquad \frac{F_P^*}{F_C^*}\longrightarrow0.
\]
Every common fixed fee $F\in(F_P^*,F_C^*)$ makes joint settlement
strictly preferable in $C$ and sequential settlement strictly preferable
in $P$, after both institutions choose their allocations. If
$M_P\geq M_C$ coordinatewise, this comparison is generated by a
productive improvement holding the other primitives fixed.
For any fee $F>0$ held fixed as $\tau\downarrow0$, sequential settlement
instead eventually dominates in both economies.
\end{corollary}

\begin{proof}[Proof sketch]
Subtracting $\tau$ does not change either asymptotic order in
Theorem~\ref{thm:optimizedloss}. Their ratio tends to zero, so the
fee interval is nonempty. Apply Proposition~\ref{prop:feethreshold}
to each economy. Both thresholds tend to zero, giving the last claim.
\end{proof}

The threshold makes the institutional choice operational. At a given
approval price, the authority purchases joint authorization when its
incremental resource cost is below the execution loss it avoids. Productive
releases lower that willingness to pay. The common-fee interval therefore
identifies programs for which an improvement in release capacity changes
the chosen authorization technology. As ordinary approvals become arbitrarily
cheap, both thresholds converge to zero and any fixed positive joint charge
eventually favors sequential execution.

\subsection{Standing balances and the decision to borrow capacity}
\label{sec:standingcapacity}

Standing authorization capacity changes the cost of execution. To identify
this margin, give the authority original balances $s r$, where $s\geq0$
and $r\gg0$ is the reserve direction. Original capacity may be consumed;
only the additional reserve $\epsilon r$ must be returned. Zero rented
reserve is admissible whenever the original ledger supports a schedule. Write $V(s)$
for the planner value and $L(s;\rho,\tau)$ for the optimized sequential
loss relative to that value. The technology has a blocked start at
$s=0$. For $s>0$, the authority can use its standing balances immediately.

The next result uses a condition on planner allocations over an interval
$0\leq s\leq\bar s$. A fixed voluntary core $d=\sum_{j=1}^k u^j$
has nonnegative release $p=Md$. For each $s$, a planner optimum can be
written
\[
 z_s^*=\alpha_s d+\sum_{j=1}^{\ell_s}v_s^j,
 \qquad 0<\underline\alpha\leq\alpha_s\leq\overline\alpha<\infty,
 \qquad \ell_s\leq\ell,
\]
where the completion bundles are voluntary and satisfy
\[
 sr+\alpha_s p+M\sum_{j=1}^t v_s^j\geq0
 \qquad(0\leq t\leq\ell_s).
\]
We call this a uniform core and completion certificate. It describes
technologies in which standing capacity changes the amount assigned to
a common executable core and finances a bounded number of terminal
adjustments. It is a finite condition on activity bundles and account
balances. When the bundles and $\alpha_s$ are affine in $s$, feasibility
can be checked at the two endpoints of the interval; a common planner
dual solution verifies optimality throughout it. Appendix~F%<xref:app:standingcapacity>

states this finite verification precisely.

\begin{theorem}[Implementation losses with standing balances]
\label{thm:standingcapacity}
Retain the positive-linear, complete-access assumptions of
Section~\ref{sec:generalproductive}, a blocked start at zero balances,
$V(0)>0$, and a uniform core and completion certificate. Fix $\rho>0$.
Uniformly over $s\in[0,\bar s]$, as $\tau\downarrow0$:
\begin{enumerate}
 \item If the technology has a strictly positive nonexpansion
 certificate, then
 \[
 L_C(s;\rho,\tau)\asymp
 \Phi_C(s;\rho,\tau)
 :=\inf_{\epsilon\geq0}
 \left\{\rho\epsilon+\frac{\tau}{s+\epsilon}\right\}
 \asymp\frac{\tau}{s+\sqrt{\tau/\rho}}.
 \]
 \item If the core has $p\ne0$ and satisfies
 $M\sum_{j=1}^t u^j\geq-\kappa p$ for every prefix and some
 $\kappa>0$, then
 \[
 \begin{aligned}
 L_P(s;\rho,\tau)&\asymp\Phi_P(s;\rho,\tau)\\
 &:=\inf_{\epsilon\geq0}
 \left\{\rho\epsilon+
       \tau\log\left(1+\frac1{s+\epsilon}\right)\right\}\\
 &\asymp\tau\log\left(1+\frac1{s+\tau/\rho}\right).
 \end{aligned}
 \]
\end{enumerate}
Here $\asymp$ denotes two-sided bounds with positive constants independent
of $s$ and sufficiently small $\tau$. The constants may depend on the
fixed economy, the certificates, $\bar s$, and $\rho$.
\end{theorem}

\begin{proof}[Proof sketch]
The initial usable stock is $(s+\epsilon)r$. The nonexpansion and
blocked-start arguments bound approvals from below using this total
stock. The bounds apply to every allocation with surplus at least
$V(0)/2$; choosing a lower-surplus allocation incurs a fixed welfare
loss. For the upper bounds, execute the core using total stock
$s+\epsilon$, by divided execution under conservation and geometric
expansion under productivity. The completion inequalities preserve
the entire rented reserve. These constructions attain the planner
optimum, so their losses contain no allocation approximation error.
Minimization over $\epsilon$ gives the displayed envelopes.
Appendix~F%<xref:app:standingcapacity>
 supplies the uniform bounds.
\end{proof}

Put $\eta=\tau/\rho$. The conserving envelope is exactly
\[
 \Phi_C=
 \begin{cases}
 2\sqrt{\rho\tau}-\rho s,&s\leq\sqrt\eta,\\
 \tau/s,&s\geq\sqrt\eta.
 \end{cases}
\]
The productive envelope is minimized at total capacity
$\max\{s,t_P\}$, where $t_P=(\sqrt{1+4\eta}-1)/2\sim\eta$.
These are minimizers of the bounding envelopes. The following
implementation decisions follow from the matching bounds themselves.

\begin{table}[!htbp]
\centering
\caption{Standing capacity and optimized implementation loss}
\label{tab:standingcapacity}
\begin{tabularx}{\linewidth}{@{}>{\raggedright\arraybackslash}p{0.23\linewidth}>{\raggedright\arraybackslash}X>{\raggedright\arraybackslash}X@{}}
\toprule
Standing capacity & Conserving technology & Productive technology\\
\midrule
$s\lesssim\eta$
& Loss $\asymp\sqrt{\rho\tau}$; borrowing of order $\sqrt\eta$ attains this rate.
& Loss $\asymp\tau\log(1/\eta)$; borrowing of order $\eta$ attains this rate.\\[3pt]
$\eta\ll s\ll\sqrt\eta$
& Loss $\asymp\sqrt{\rho\tau}$; any implementation attaining this order must borrow capacity of order $\sqrt\eta$.
& Loss $\asymp\tau\log(1/s)$; existing balances suffice to attain this order.\\[3pt]
$\sqrt\eta\lesssim s\to0$
& Loss $\asymp\tau/s$; existing balances suffice.
& Loss $\asymp\tau\log(1/s)$; existing balances suffice.\\[3pt]
Fixed $s>0$
& Loss $\asymp\tau$, with a coefficient of order $1/s$ as $s$ becomes small.
& Loss $\asymp\tau$, with a coefficient of order $\log(1/s)$ as $s$ becomes small.\\
\bottomrule
\end{tabularx}
\end{table}

The middle regime identifies a financing decision: productive releases
allow the authority to implement its planner allocation using standing
capacity, whereas conservation requires additional authorization to
attain the optimal welfare-loss order. At fixed positive balances,
both technologies' sequential losses become linear in the approval price.
Their coefficients retain the economically relevant distinction.

\begin{corollary}[Institutional choice with positive standing balances]
\label{cor:standingreversal}
Consider a conserving economy and a productive economy satisfying the
respective conditions of Theorem~\ref{thm:standingcapacity}, with common
preferences, action caps, trading permissions, $s,\rho,\tau$, and reserve direction. For every path
$s=s(\tau)\to0$, including strictly positive paths,
\[
 L_P(s;\rho,\tau)/L_C(s;\rho,\tau)\longrightarrow0.
\]
Moreover, for every sufficiently small fixed $s>0$, there is a nonempty
interval of constants $f_s>0$ such that the common joint commitment
fee $F=f_s\tau$ makes joint settlement optimal in the conserving
economy and sequential settlement optimal in the productive economy,
for all sufficiently small $\tau$. Each institution chooses its
allocation. The same conclusion applies to a technological improvement
when preferences, caps, and trading permissions are common and
the two technologies additionally satisfy $M_C\leq M_P$.
\end{corollary}

\begin{proof}[Proof sketch]
The envelopes imply the ratio limit. For fixed $s$, uniform constants
give $L_C\geq c_C\tau/s$ and
$L_P\leq C_P\tau\log(1+1/s)$ for small $\tau$.
As $s\downarrow0$, the first coefficient exceeds the second.
Choose
\[
 \max\{0,C_P\log(1+1/s)-1\}<f_s<c_C/s-1
\]
and apply the threshold $F^*=L-\tau$. Planner values remain bounded
away from zero, ensuring adoption. The full proof is in
Appendix~F%<xref:app:standingcapacity>
.
\end{proof}


\section{Exact costs, endogenous allocations, and robustness}
\label{sec:optimalreserve}
\subsection{The exact cost of a proportional target}
\label{sec:productivecost}

For $0\leq h\leq1/4$, put
\[
 Q_h(\theta;\rho,\tau)
 =\inf_{\epsilon>0}\{\rho\epsilon+\tau N_h(\theta,\epsilon)\}.
\]
This is a cost conditional on implementing $z_\theta$. The endogenous welfare
comparison below allows each institution to choose any feasible allocation.

\begin{theorem}[The exact cost frontier]
\label{thm:productivecost}
For $0\leq h\leq1/4$, $\theta\in(0,1]$, and $\rho,\tau>0$,
\begin{equation}
 Q_h(\theta;\rho,\tau)=\min_{T\in\{2,3,\ldots\}}
    \{\rho\epsilon_T(h,\theta)+\tau T\}.
 \label{eq:productivecostexact}
\end{equation}
An optimal reserve and integer approval count exist. The reserve infimum can
equivalently be taken over $0<\epsilon\leq\theta/4$.
For each fixed $h>0$ and $\theta>0$, as $\tau/\rho\downarrow0$,
$Q_h=\Theta(\tau\log(\rho/\tau))$, whereas
$Q_0=\Theta(\sqrt{\rho\theta\tau})$.
\end{theorem}
\begin{proof}[Proof sketch]
At each approval count the exact frontier gives the least reserve; minimizing
over integers proves the formula. Its objective tends to infinity with $T$,
so a minimum exists. Two equal pair increments implement the target from
reserve $\theta/4$, making larger reserves unnecessary. At $h=0$, minimize
$\rho\theta/[12(T-1)]+\tau T$. At fixed $h>0$, the exact reserve decreases
geometrically in $T$; a count proportional to $\log(\rho/\tau)$ attains
the upper order, and the geometric lower bound gives the matching order.
Appendix~B%<xref:app:localrobustness>
 supplies the full argument.
\end{proof}

\subsection{Reversal when each institution chooses its allocation}
\label{sec:endogenousreversal}

Let $\mathcal V_h^J$ and $\mathcal V_h^{\mathrm{seq}}$ denote the best joint
and sequential employee surplus net of implementation costs, including the
zero-cost baseline. The joint value is
\[
 \mathcal V_h^J=\max\{0,V_h^*-F-\tau\}.
\]
The sequential value is the supremum of
$W(z)-\rho\epsilon-\tau T$ over all voluntary admissible schedules, all
$\epsilon>0$, and all terminal allocations whose reserve can be fully returned,
also including zero. Sequential settlement can select every feasible endpoint and its
admissible compensated schedule.

\begin{theorem}[Endogenous-allocation institutional reversal]
\label{thm:endogenousreversal}
In the family $M_h$, fix $\rho=12$, $\tau=1/10000$, $\mathcal K=0$, and charge
$F=1/100$ for every nonzero joint package. Then
\[
 \mathcal V_0^J=\frac{8899}{10000}
 >\frac{8799}{10000}=\mathcal V_0^{\mathrm{seq}},
\]
whereas
\[
 \mathcal V_{1/4}^{\mathrm{seq}}
 \geq\frac{27280647}{30064375}
 >\frac{98889}{110000}=\mathcal V_{1/4}^J.
\]
Thus, under the same fixed implementation costs, productive releases reverse
the welfare ranking of joint and sequential settlement even when each
institution may choose its allocation. The productive lower bound is attained
by a voluntary 16-approval schedule implementing $z_{1/4}^*$ with reserve
$1/144309$ in every account and full repayment.
\end{theorem}

\begin{proof}[Proof sketch]
At $h=0$, repayment and conservation imply that every feasible terminal
allocation is $z_\theta$, $0\leq\theta\leq1$. For $T\geq2$, the exact
frontier makes net surplus at most
\[
 \frac9{10}\theta-\frac{\theta}{T-1}-\frac{T}{10000}.
\]
For each $T$, this is linear in $\theta$, so a positive optimum uses
$\theta=1$. Writing $n=T-1$, the minimum of
$1/n+(n+1)/10000$ is $201/10000$, attained at $n=100$.
The best joint value is $9/10-101/10000$.

At $h=1/4$, Proposition~\ref{prop:productiveplanner} gives the global planner
value $10/11$. Proposition~\ref{prop:exactproductive} implements
$z_{8/11}$ with $16$ approvals and reserve $e=1/144309$. The target releases
$1/11$ of pension capacity. A final supplier-2 approval with demander increment
$3/11$ and supplier increment $2/11$ consumes that release, adds $1/11$ of
position capacity, and reaches $(1,4/11,6/11)$. It has positive surplus and
leaves the reserve repayable. Its implementation cost is
\[
 \overline C=\frac{12}{144309}+\frac{16}{10000}
            =\frac{50603}{30064375}<\frac{101}{10000}.
\]
The last first-stage approval and the completion use the same pair.
Combining them preserves linear feasibility and balanced voluntary compensation,
so only 16 approvals are charged. Hence its net value exceeds the best joint value, whose allocation is also
chosen optimally. Appendix~G%<xref:app:endogenousproof>
 verifies every approval,
transfer, cap, and account balance. The productive sequential number is an
attainable lower bound, which suffices for the optimized ranking.
\end{proof}

\begin{table}[!htbp]
\centering
\caption{Endogenous allocations and the reversal of institutional welfare}
\label{tab:endogenousreversal}\label{tab:productivereversal}
\begin{tabularx}{\linewidth}{@{}>{\raggedright\arraybackslash}Xrrl@{}}
\toprule
Technology & Best joint & Best sequential & Higher value\\
\midrule
Conserving ($h=0$) & $0.889900$ & $0.879900$ & Joint\\
Productive ($h=1/4$) & $0.898991$ & $\geq\underline V_P$ & Sequential\\
\bottomrule
\end{tabularx}

\smallskip
\begin{minipage}{\linewidth}
The productive lower bound is
$\underline V_P=27280647/30064375\simeq0.907408$.
Decimals are rounded; the inequality refers to the exact fraction.
Both joint entries and the conserving sequential entry are global optimized
values. The productive sequential entry uses a feasible 16-approval schedule.
Preferences, caps, prices, and the fixed fee on every nonzero joint package
are unchanged.
\end{minipage}
\end{table}

The exact productive frontier and cost theorem concern proportional targets.
They provide the first stage of this proof, while the endogenous planner
calculation supplies the correct endpoint. Theorem~\ref{thm:endogenousreversal}
therefore strengthens the fixed-target comparison in the online appendix
by implementing the endogenous planner optimum. The appendix displays the 17-stage construction and
the merged final approval explicitly. The feasible witness establishes the productive lower bound required for
the ranking; both joint values and the conserving sequential value are exact.

Corollary~A.3%<xref:cor:productivereversal>
 in the online appendix gives the
corresponding fixed-target comparison and its exact cost-minimizing reserves.

\subsection{When standing capacity changes the institution}\label{sec:standingdecisions}
The three-account family supplies an explicit certificate. For
$0\leq h<1/2$ and $0\leq s\leq1/32$, original balances $s\mathbf1$
give
\[
 z_{h,s}^*=\left(1,\frac{1-8s}{2+3h},
 \frac{1+2h+(4+2h)s}{2+3h}\right),\qquad
 V_{h,s}^*=1-\frac{1+h}{5(2+3h)}
                 +\frac{2-h}{5(2+3h)}s.
\]
Implement the proportional target with
$\theta_s=2(1-8s)/(2+3h)$ using total initial capacity
$s+\epsilon$. One supplier-2 approval then adds
\[
 (\Delta a,\Delta b_2)=
 \left(\frac{3h+16s}{2+3h},
       \frac{2h+(12+2h)s}{2+3h}\right).
\]
It reaches $z_{h,s}^*$, consumes the permitted original balances,
and returns $\epsilon$ in each account. Compensation equals
$\Delta b_2/10$, leaving the supplier indifferent and the demander
strictly better off whenever the completion is nonzero.

\begin{proposition}[An exact frontier with standing capacity]
\label{prop:standingexact}
In $M_0$ with original balances $s\mathbf1$, $0\leq s\leq1/32$,
the minimum additional reserve implementing the planner optimum in
at most an integer $T\geq2$ voluntary approvals is
\[
 \epsilon_T^*(s)=
 \left[\frac{1-8s}{12(T-1)}-s\right]_+.
\]
In particular, for $s>0$ the minimum number of approvals implementing
that optimum without rented capacity is
\[
 N_0(s)=\left\lceil1+\frac{1-8s}{12s}\right\rceil.
\]
\end{proposition}

\begin{proof}[Proof sketch]
Let $t=s+\epsilon$ and let $q$ be the difference between payroll
and position balances. Conservation bounds $|q|\leq3t$ at every
prefix. Initially $q=0$; at the planner endpoint $q=-3s$.
Each supplier's activity equals twice its contribution to the
variation of $q$. Hence total supplier activity satisfies
$1-2s\leq12t(T-1)+6s$, giving the lower bound.
Alternating boundary transactions attain this variation, with a
final supplier-2 transaction that leaves the full reserve available
for repayment. When standing balances exceed the required amplitude,
smaller alternating excursions attain the same endpoint without
renting capacity. Every transaction has positive employee surplus.
Appendix~F.5%<xref:app:standingexact>
 gives the construction and proves
optimality over all schedules.
\end{proof}

The formula separates two uses of original capacity. A larger stock
directly replaces rented capacity, and it changes the planner's
supplier mix, reducing the quantity that must pass through repeated
conversions. Because approvals are discrete, institutional choice
can include an interval of indifference rather than a single cutoff.

\begin{corollary}[Standing capacity determines when productivity changes the institution]
\label{cor:standingcutoffs}
Fix $\rho=12$, $\tau=1/10000$, $F=1/100$, and original balances
$s\mathbf1$, with $0\leq s\leq1/32$. Define
\[
 s_J=\frac1{1208},\qquad
 s_S=\frac{9901}{11960000}.
\]
After optimizing allocations, reserves, and schedules, the
conserving economy $M_0$ strictly prefers joint settlement for
$s<s_J$, is indifferent for $s\in[s_J,s_S]$, and strictly prefers
sequential settlement for $s>s_S$. The productive economy
$M_{1/4}$ strictly prefers sequential settlement throughout the
interval $[0,1/32]$.
\end{corollary}

\begin{proof}[Proof sketch]
For $n=T-1\geq2$, the conserving sequential optimum conditional
on $n$ uses the planner allocation and has implementation cost
\[
 C_n(s)=\left[\frac{1-8s}{n}-12s\right]_+
          +\frac{n+1}{10000}.
\]
Schedules with one or two approvals, including their optimal
quantity choices, cannot improve on joint settlement below the
stated cutoff. The $n=100$ schedule first matches the joint cost
$101/10000$ at $s_J$; the $n=99$ schedule first matches it
at $s_S$ and falls below it for larger stocks. Every other approval count crosses later or has an
approval cost already exceeding the joint cost. In $M_{1/4}$,
the proportional construction and one completion reach the new
planner optimum using at most 17 approvals and rented reserve
at most $1/144309$. Their cost is strictly below $101/10000$
for every stated $s$. Appendix~F.5%<xref:app:standingexact>
 verifies
the global optimization and the cutoff inequalities.
\end{proof}

This comparison holds prices and commitment costs fixed while
varying standing authorization capacity. Productivity changes the
chosen institution when $s<s_J$. Once $s>s_S$, existing capacity
already makes sequential settlement preferable under conservation,
so productivity improves its performance without changing the
institution. The interval between these cutoffs reflects the
integer number of approvals.

For the same prices as Theorem~\ref{thm:endogenousreversal}, take
$s=1/2500$. In the productive economy, nine proportional approvals
and the completion implement the planner optimum using only original
balances. Table~\ref{tab:positivereversal} reports exact joint values,
the exact conserving optimum, and this attainable productive value.
The fixed fee is $F=0.01$ in both economies. Joint productive value is
rounded to six decimals; its exact value is $494473/550000$. The
productive sequential lower bound is rounded downward from
$249739/275000$. The two conserving entries are exact.

\begin{table}[!htbp]
\centering
\caption{Optimized institutional comparison with positive original balances}
\label{tab:positivereversal}
\begin{tabularx}{\linewidth}{@{}>{\raggedright\arraybackslash}Xrrl@{}}
\toprule
Technology & Best joint value & Best sequential value & Preferred rule\\
\midrule
$h=0$ & $0.889980$ & $0.884812$ & Joint\\
$h=1/4$ & $0.899042$ & $\geq0.908141$ & Sequential\\
\bottomrule
\end{tabularx}
\end{table}

All accounts initially contain positive capacity. Productivity changes
the authority's institution and removes the need to rent capacity in
the displayed planner implementation. Appendix~F%<xref:app:standingcapacity>

verifies the fractions, compensation, intermediate balances, and global
optimization; Appendix~F.5%<xref:app:standingexact>
 proves the exact
conserving value.


\paragraph{The conserving optimum.}
Proposition~A.4%<xref:prop:exactdesign>
 gives the full conserving solution.
With common verification cost $\mathcal K(z_\theta)=k\theta$, the optimal
positive quantity is one. The authority compares the joint fee with
$G(\rho,\tau)=\min_{n\geq1}\{\rho/(12n)+\tau n\}$.
The minimizing integer satisfies
$\tau n(n-1)\leq\rho/12\leq\tau n(n+1)$; the sequential schedule
uses $n+1$ approvals and reserve $1/(12n)$. This exact choice connects
the resource frontier to the inventory balance in \citet{Baumol1952}.

\paragraph{Allocation and participation margins.}
The comparison also permits the joint fee to vary with the allocation.
Proposition~A.5%<xref:prop:allocationfees>
 proves that every common fee profile
in $[0.005,0.015]$ preserves the reversal; a coverage charge can change the
joint quantity to $0.99z_h^*$. Proposition~A.6%<xref:prop:allocationdistortion>

shows that concave benefits can reduce the sequentially chosen quantity to
$7/8$, and that employee participation can raise the two-approval reserve
above its physical minimum. The appendix gives the complete institutional
optimizations and compensated schedules.

\paragraph{Fixed adjustment units and original slack.}
The reversal also holds with a fixed adjustment unit and positive
original account balances. The following extension imposes the same activity grid
on both institutions; the core model remains divisible. Let
$\mathcal V^J_{h,\delta}(R^0)$ and
$\mathcal V^{\mathrm{seq}}_{h,\delta}(R^0)$ be their optimized net values
when every activity increment is a multiple of $\delta$, original
balances are $R^0$, and every rented reserve is returned. Transfers and
reserve capacity remain divisible, and the authority selects the
allocation and schedule.

\begin{corollary}[Reversal with fixed units and original account slack]
\label{cor:gridrobustreversal}
Maintain the prices and financing assumptions of
Theorem~\ref{thm:endogenousreversal}. Fix $\delta=1/500$ and let
$R^0\in[0,1/2500]^3$. Then
\[
 \mathcal V^J_{0,\delta}(R^0)=\frac{8899}{10000}
 >\frac{44399}{50000}
 \geq\mathcal V^{\mathrm{seq}}_{0,\delta}(R^0),
\]
whereas
\[
 \mathcal V^{\mathrm{seq}}_{1/4,\delta}(R^0)
 \geq\frac{9013}{10000}
 >\frac{8989}{10000}
 =\mathcal V^J_{1/4,\delta}(R^0).
\]
The productive lower bound uses a 17-approval schedule with reserve
$1/2000$ per account, implements a grid-constrained planner optimum,
and returns that reserve in full. The result includes every strictly
positive original ledger in the displayed box.
\end{corollary}

\noindent\emph{Proof sketch.}
An exact grid schedule implements $(1,183/500,272/500)$, with gross
surplus $909/1000$. Its reserve and approval costs total $77/10000$.
Feasibility is preserved when original balances increase. In the
conserving economy, a variation bound on the difference between payroll
and position balances gives the global upper bound
$\mathcal V^{\mathrm{seq}}_0(s\mathbf 1)
\leq8799/10000+(101/5)s$, even with divisible activities.
Planner dual bounds and the activity grid give the two exact joint
values. Monotonicity in the original ledger extends the bounds to
the entire box. Appendix~J%<xref:app:gridrobustreversal>
 supplies the
schedule and proof. The specified grid and slack box therefore preserve the strict
institutional ranking at the common implementation prices.

\section{An empirical application: priority in water allocation trade}
\label{sec:waterapplication}

Administrative records from the Victorian Water Register show how processing
order changes the allocation of a positive stock of authorized trade
opportunity. The application measures successful requests and executed water
volume under alternative priority rules. These outcomes make the authority's
allocation decision observable in an institution with binding capacity.
Requests are treated as indivisible at their submitted volumes, connecting
the exercise to Appendix~L%<xref:app:indivisibility>
.

\subsection{Institution and administrative replay}
For the October 2025 Goulburn to Murray opening, the authority randomized
applications submitted during a common window. The Australian Competition
and Consumer Commission independently observed the randomization
\citep{ACCC2025October}. The published application table contains 203 eligible
requests, submitted on 15 October and processed on 16 October. Each row
records the requested volume, processing order, outcome, and trade opportunity
before and after processing. Receipt identifiers also recover arrival order
\citep[Table~5 and notes]{DEECA2025October}.

The stock available immediately before this queue was
$R=31{,}215.1$ megalitres (ML). It equals the opening opportunity of
$31{,}523.4$ ML, less $326.3$ ML of earlier tagged use, plus $18.0$ ML of
reverse trade \citep{DEECA2025October}. The reverse trade replenished usable
corridor capacity before the queue began. It therefore enters the starting
stock held fixed across rules. No tagged use occurred during the queue.
The common pool contains the 203 eligible applications; it excludes two
ineligible submissions and a later trade that used the final $5.3$ ML.

Let $v_i>0$ be the submitted volume and $q$ a permutation of the pool
$\mathcal I$. The capacity benchmark accepts a request in full whenever it
fits, and continues after a refusal:
\[
 Y_i(q)=\mathbf 1\!\left\{
 v_i\le R-\sum_{j\prec_q i}v_jY_j(q)\right\},\qquad
 A(q)=\sum_iY_i(q),\qquad Q(q)=\sum_i v_iY_i(q).
\]
Here $A$ counts successful applications and $Q$ measures executed volume.
Exact arithmetic reproduces every observed October decision and every
intermediate balance. Appendix~M%<xref:app:waterbenchmark>
 documents the
391 records covering October and December, the exclusions, and the replay.

\subsection{Priority, application counts, and capacity use}
Write $v_{(1)}\le\cdots\le v_{(203)}$ for the ordered request sizes.
The maximum application count under the capacity benchmark is
\[
 A^*(R)=\max\left\{k:\sum_{j=1}^{k}v_{(j)}\le R\right\}.
\]
Processing the smallest requests first attains this bound: every subset of
$k$ requests consumes at least the sum of the $k$ smallest volumes.
The maximum need not select a unique order. In these data,
\[
 \sum_{j=1}^{159}v_{(j)}=31{,}193.0
 \ \le\ 31{,}215.1
 \ <\ 31{,}693.0=\sum_{j=1}^{160}v_{(j)}.
\]
Thus $159$ is the exact maximum, and executing $160$ requests would require
at least another $477.9$ ML of trade opportunity under this benchmark.

\begin{table}[htbp]
\centering\small
\caption{October water trade: the same applications and starting capacity}
\label{tab:waterqueues}
\begin{tabular}{@{}lrrr@{}}
\toprule
Processing rule & Successful applications & Volume (ML) & Unused (ML)\\
\midrule
Recorded random priority & 98 & 31,209.8 & 5.3\\
Arrival order & 106 & 31,206.4 & 8.7\\
Smallest volume first & 159 & 31,193.0 & 22.1\\
\bottomrule
\end{tabular}
\par\smallskip
\begin{minipage}{\linewidth}\footnotesize
The first row reproduces the administrative record. The other rows are
capacity counterfactuals for the same 203 submitted requests, assuming all
other requirements remain satisfied. All rules process 203 applications.
Smallest first breaks equal volumes by recorded priority. Source:
\citet[Table~5]{DEECA2025October}; calculations in
Appendix~M%<xref:app:waterbenchmark>
.
\end{minipage}
\end{table}

The change from recorded priority to smallest first yields 61 more successful
applications, a $62.2\%$ increase relative to 98. The gap is $38.4\%$ of
the attainable count. With the stated tie rule, 79 observed successes remain
selected, 19 are displaced, and 80 previously refused requests enter.
Executed volume falls by only $16.8$ ML. Near exhaustion of the volume limit
therefore coexists with a substantial difference in how many applications
receive capacity. A simulation of 100,000 uniform permutations gives a mean
of $93.66$ successes and a central $95\%$ range of 75 to 112. The recorded
98 successes are typical of this conditional lottery benchmark.

\subsection{Economic interpretation and account requirements}
The application identifies an allocation margin: priority determines which
requests use existing capacity. A count objective places equal weight on
applications. A welfare objective instead requires their gains from trade;
191 of the 203 requests are classified as related party trades, and 174
record a zero price. These entries do not identify those gains. Applications
also need not represent distinct traders. A permanent change in priority
could change submission sizes and application splitting, whereas the reported
comparison holds the submitted menu fixed.

The comparison clarifies the role of settlement in the theory. In the scalar
capacity benchmark, any feasible selected subset can be executed sequentially
in any order. Joint selection and an optimal sequential rule consequently
attain the same count. The measured gap concerns the particular priority
rule. The main theorems additionally characterize the cost of execution when
separate accounts and productive releases constrain the authority's chosen
allocation. Moreover, the 203 administrative decisions are distinct from
the successful application counts in Table~\ref{tab:waterqueues}; those
counts cannot be substituted for the costly approval count in the theory.

December provides a direct observation of an additional account requirement.
At priority 45, an application for $62.3$ ML was refused with $23{,}824.0$ ML
of trade opportunity remaining. The report attributes the refusal to an
insufficient seller balance \citep[Table~5, note~8]{DEECA2025December}.
The seller account is a private entitlement alongside the public trade
limit. Its unobserved links and balance history determine whether a different
schedule could fund that seller. The October count maximum is therefore
conditional on the other account requirements. Together, the records
establish a substantial allocation margin and distinct execution constraints.

The $18.0$ ML reverse trade illustrates how completed transactions replenish
capacity for later use. This addition enters the initial stock $R$ held
fixed in our comparisons. The productive-release theorem requires a
voluntary sequence that finances its intermediate steps, supports repeated
expansion, and permits completion at the chosen planner allocation. The
records establish an operational channel for replenishment; the theorem
identifies when that channel changes optimized implementation costs and
institutional choice.

\section{Conclusion}
\label{sec:conclusion}

Productive releases change the cost of implementing a planner-optimal
allocation under separate public accounts. Conserving transactions repeatedly
use existing authorization capacity, whereas an executable productive sequence
creates capacity for subsequent approvals. The exact frontier and the general
loss bounds quantify this distinction after the authority chooses both its
allocation and its schedule.

Standing balances determine how the mechanism enters institutional decisions.
With empty accounts, conservation and productivity generate different
small-cost rates. With fixed positive balances, both losses are linear in the
approval price, but the required approval counts depend inversely and
logarithmically, respectively, on scarce standing capacity. The uniform
comparison identifies a range in which productive implementation attains the
minimum-loss order using existing balances, while conserving implementation
gains from temporary reserves. Productive releases therefore reduce the
standing capacity needed to economize on paid authorization and joint
commitment. The commitment thresholds translate this quantitative advantage
into a reversal of the preferred settlement rule, with allocations endogenous
under both institutions. The exact conserving frontier identifies a
standing-capacity cutoff above which both technologies select sequential
settlement at the same fixed implementation prices.

The water trade application makes the allocation margin concrete. With
positive initial trading opportunity and fixed submitted requests, different
priority rules produce nearly identical volume use and sharply different
application counts. Tracking the selected allocation and the underlying
account constraints is therefore essential to evaluating an execution rule.
The theory supplies the further comparison when productive releases and
separate public accounts determine the cost of implementing that allocation.

The framework preserves the distinction between employee compensation and
public authorization. Transfers secure participation; resource releases
finance execution. The main results retain complete information and the
$B\equiv0$ mandate. The appendix extends the analysis to valued public
balances and discrete decisions, identifying the financing conditions and
reserve floors relevant to those environments. These results provide a basis
for comparing account reforms through their effects on standing capacity,
temporary authorization, and approval requirements.


\begin{thebibliography}{99}
\bibitem[Arrow and Debreu(1954)]{ArrowDebreu1954}
\textsc{Arrow, Kenneth J., and Gerard Debreu} (1954): ``Existence of an equilibrium for a competitive economy,'' \emph{Econometrica}, 22(3), 265--290. DOI: \url{10.2307/1907353}.

\bibitem[Atalay et~al.(2010)]{AtalayMartinMcAndrews2010}
\textsc{Atalay, Enghin, Antoine Martin, and James McAndrews} (2010):
``The welfare effects of a liquidity-saving mechanism,'' Staff Report 331,
Federal Reserve Bank of New York, June 2008, revised January 2010.
\url{https://www.newyorkfed.org/research/staff_reports/sr331.html}.


\bibitem[ACCC(2025)]{ACCC2025October}
\textsc{Australian Competition and Consumer Commission} (2025):
``Goulburn to Murray trade opening, 15 October 2025: ACCC independent observer
of the execution of randomisation tool,'' statement, 24 October.
\url{https://waterregister.vic.gov.au/images/documents/ACCC_statement_randomisation_october_2025.pdf}.

\bibitem[Baumol(1952)]{Baumol1952}
\textsc{Baumol, William J.} (1952): ``The transactions demand for cash: An inventory theoretic approach,'' \emph{The Quarterly Journal of Economics}, 66(4), 545--556. DOI: \url{10.2307/1882104}.

\bibitem[Boyd and Vandenberghe(2004)]{BoydVandenberghe2004}
\textsc{Boyd, Stephen, and Lieven Vandenberghe} (2004): \emph{Convex Optimization}. Cambridge University Press. DOI: \url{10.1017/CBO9780511804441}.

\bibitem[CPSS and IOSCO(2012)]{CPSSIOSCO2012}
\textsc{Committee on Payment and Settlement Systems and Technical Committee of the International Organization of Securities Commissions} (2012):
\emph{Principles for Financial Market Infrastructures}. Basel: Bank for International Settlements, April.
\url{https://www.bis.org/publications/principles-financial-market-infrastructures.pdf}.

\bibitem[DEECA(2025a)]{DEECA2025October}
\textsc{Department of Energy, Environment and Climate Action} (2025a):
\emph{Goulburn-to-Murray Post Trade Opening Report: October 2025}.
Victorian Water Register, 23 October.
\url{https://waterregister.vic.gov.au/images/documents/Goulburn-to-Murray-Post-Trade-Opening-Report-October-2025.pdf}.

\bibitem[DEECA(2025b)]{DEECA2025December}
\textsc{Department of Energy, Environment and Climate Action} (2025b):
\emph{Goulburn-to-Murray Post Trade Opening Report: December 2025}.
Victorian Water Register, 19 December.
\url{https://waterregister.vic.gov.au/images/documents/Goulburn-to-Murray-Post-Trade-Opening-Report-December-2025.pdf}.

\bibitem[Feldman(1973)]{Feldman1973}
\textsc{Feldman, Allan M.} (1973): ``Bilateral trading processes, pairwise optimality, and Pareto optimality,'' \emph{Review of Economic Studies}, 40(4), 463--473. DOI: \url{10.2307/2296581}.

\bibitem[Fisher(1989)]{Fisher1989}
\textsc{Fisher, Franklin M.} (1989): ``It takes t$^{*}$ to tango: Trading coalitions with fixed prices,'' \emph{Review of Economic Studies}, 56(3), 391--403. DOI: \url{10.2307/2297554}.

\bibitem[Fraca and Haddad(2015)]{FracaHaddad2015}
\textsc{Fraca, Estíbaliz, and Serge Haddad} (2015): ``Complexity analysis of continuous Petri nets,'' \emph{Fundamenta Informaticae}, 137(1), 1--28. DOI: \url{10.3233/FI-2015-1168}.

\bibitem[Goldman and Starr(1982)]{GoldmanStarr1982}
\textsc{Goldman, Steven M., and Ross M. Starr} (1982): ``Pairwise, t-wise, and Pareto optimalities,'' \emph{Econometrica}, 50(3), 593--606. DOI: \url{10.2307/1912603}.

\bibitem[J{\'u}lvez et~al.(2003)]{JulvezRecaldeSilva2003}
\textsc{Júlvez, Jorge, Laura Recalde, and Manuel Silva} (2003): ``On reachability in autonomous continuous Petri net systems,'' in \emph{Applications and Theory of Petri Nets 2003}, ed. by Wil M. P. van der Aalst and Eike Best. \emph{Lecture Notes in Computer Science}, 2679. Springer, 221--240. DOI: \url{10.1007/3-540-44919-1_16}.

\bibitem[Kahn and Roberds(2001)]{KahnRoberds2001}
\textsc{Kahn, Charles M., and William Roberds} (2001): ``Real-time gross settlement and the costs of immediacy,'' \emph{Journal of Monetary Economics}, 47(2), 299--319. DOI: \url{10.1016/S0304-3932(01)00047-2}.

\bibitem[Kahn et~al.(2003)]{KahnMcAndrewsRoberds2003}
\textsc{Kahn, Charles M., James McAndrews, and William Roberds} (2003): ``Settlement risk under gross and net settlement,'' \emph{Journal of Money, Credit and Banking}, 35(4), 591--608.
\url{https://www.newyorkfed.org/research/staff_reports/sr86.html}.

\bibitem[Koeppl et~al.(2008)]{KoepplMonnetTemzelides2008}
\textsc{Koeppl, Thorsten V., Cyril Monnet, and Ted Temzelides} (2008): ``A dynamic model of settlement,'' \emph{Journal of Economic Theory}, 142(1), 233--246. DOI: \url{10.1016/j.jet.2006.10.005}.

\bibitem[Madden(1978)]{Madden1978}
\textsc{Madden, Paul} (1978): ``Why the Edgeworth process assumption isn’t that bad,'' \emph{Review of Economic Studies}, 45(2), 279--283. DOI: \url{10.2307/2297343}.

\bibitem[Martin and McAndrews(2008)]{MartinMcAndrews2008}
\textsc{Martin, Antoine, and James McAndrews} (2008): ``Liquidity-saving mechanisms,'' \emph{Journal of Monetary Economics}, 55(3), 554--567. DOI: \url{10.1016/j.jmoneco.2007.12.011}.

\end{thebibliography}
\end{document}
